% =====================================================================
%  Seance 3 -- D. TP : detecter les incidents avant la panne (frames 24--30)
% =====================================================================
\section{TP « Détecter les incidents avant la panne »}

\begin{frame}{Le TP : objectif et règle du jeu}
\begin{columns}[t]
\begin{column}{0.5\textwidth}
\begin{kbox}[Ce que vous produisez à 18 h]
Un \textbf{système d'alerte} sur \file{nosys\_features.csv} qui :
\begin{enumerate}
\item détecte sans étiquette (IsolationForest) ;
\item prédit « incident $\le$ 30 min » avec un modèle supervisé ;
\item respecte un \textbf{budget NOC} : $\le$ 5 alertes/jour/serveur ;
\item est \textbf{justifié par les chiffres} : matrice de confusion,
courbe PR, tableau comparatif.
\end{enumerate}
\end{kbox}
\begin{ebox}[Rappels de cadrage]
\begin{itemize}
\item Split \textbf{temporel} imposé : train $<$ 22/09 (24 192),
test $\ge$ 22/09 (10 368, 9 jours).
\item Graine 5534 partout (\py{random\_state=5534}).
\item Jamais de KFold mélangé.
\end{itemize}
\end{ebox}
\end{column}
\begin{column}{0.46\textwidth}
\begin{center}
\includegraphics[width=\textwidth]{f20_courbes_pr.png}
\end{center}
\vspace{1mm}
\begin{center}\small
Votre cible à battre (corrigé officiel) :\\ $F_1$ \textbf{0,769},
rappel \textbf{0,625}, AP \textbf{0,63}.
\end{center}
\end{column}
\end{columns}
\end{frame}

\begin{frame}[fragile]{TP B1 --- IsolationForest sans étiquette (25 min)}
\begin{codebox}[B1. Détecter sans jamais regarder la cible]
\begin{lstlisting}
import pandas as pd
from sklearn.ensemble import IsolationForest

f = pd.read_csv("data/nosys_features.csv", parse_dates=["fenetre"])
X = f.drop(columns=["serveur", "fenetre", "incident_30min"])
X = X.fillna(0.0)                       # cpu_delta de la 1re fenetre

train = f["fenetre"] < "2026-09-22"
iso = IsolationForest(n_estimators=300, contamination=0.01,
                      random_state=5534, n_jobs=2).fit(X[train])
f["score_anomalie"] = -iso.score_samples(X)

top20 = f.nlargest(20, "score_anomalie")
print(top20[["serveur", "fenetre", "cpu_moy", "lat_p99",
             "incident_30min"]])
\end{lstlisting}
\end{codebox}
\begin{abox}[Livrable B1]
La liste des 20 fenêtres les plus anormales, triée par score.
Contrôle attendu (corrigé) : \textbf{20/20 sont des pré-incidents}.
Notez à la main ce que vous \emph{pouviez} savoir : aucun
\py{incident\_30min} n'est entré dans le calcul.
\end{abox}
\end{frame}

\begin{frame}[fragile]{TP B2 --- recouvrement détections / alarmes réelles (20 min)}
\begin{codebox}[B2. Confronter le non-supervisé au réel]
\begin{lstlisting}
seuil = f.loc[train, "score_anomalie"].quantile(0.99)
f["anomalie"] = f["score_anomalie"] > seuil

test = f[~train]
print("anomalies test :", test["anomalie"].sum())      # 174
print("dont vrais incidents :",
      test.loc[test["anomalie"], "incident_30min"].sum())  # 40

# recouvrement : par fenetre, combien de serveurs detectes ?
cov = (test.groupby("fenetre")["anomalie"].sum()
           .value_counts().sort_index())
print(cov)     # incidents vus sur 1, 2, ... serveurs a la fois
\end{lstlisting}
\end{codebox}
\begin{columns}[t]
\begin{column}{0.48\textwidth}
\begin{wbox}[Ce que 174 anomalies veulent dire]
\textbf{40} correspondent à un incident à venir ; les
\textbf{134} autres sont des anomalies bénignes (batch, backup,
pics). Chaque fausse alerte coûte une nuit d'ingénieur : le
non-supervisé seul ne suffit pas.
\end{wbox}
\end{column}
\begin{column}{0.48\textwidth}
\begin{ebox}[Livrable B2]
Un tableau croisé (anomalie $\times$ incident) sur le test + un
commentaire de 3 lignes : quelle part des anomalies vaut le coup
d'être envoyée au NOC ?
\end{ebox}
\end{column}
\end{columns}
\end{frame}

\begin{frame}[fragile]{TP B3 --- forêt balanced et matrice de confusion (30 min)}
\begin{codebox}[B3. Entraîner, puis JUGER sur le futur]
\begin{lstlisting}
from sklearn.ensemble import RandomForestClassifier
from sklearn.metrics import confusion_matrix

Xtr, ytr = X[train], f.loc[train, "incident_30min"]
Xte, yte = X[~train], f.loc[~train, "incident_30min"]

rf = RandomForestClassifier(n_estimators=300,
        class_weight="balanced", random_state=5534,
        n_jobs=2).fit(Xtr, ytr)
proba = rf.predict_proba(Xte)[:, 1]
pred  = (proba >= 0.5).astype(int)

print(confusion_matrix(yte, pred))
\end{lstlisting}
\end{codebox}
\begin{columns}[t]
\begin{column}{0.48\textwidth}
\begin{wbox}[Le piège que vous devez rencontrer]
Avec \py{class\_weight='balanced'}, le seuil 0,50 n'est plus
calibré : au corrigé, la forêt prédit 7 positifs... tous faux
(précision 0,00). Ce n'est \textbf{pas un bug} : c'est le déplacement
du seuil effectif. Ne paniquez pas, passez au B4.
\end{wbox}
\end{column}
\begin{column}{0.48\textwidth}
\begin{ebox}[Livrable B3]
La matrice de confusion sur le test + les 4 cadres
(TN/FP/FN/TP) commentés en coût métier : qui dérange qui,
combien de SLA en jeu.
\end{ebox}
\end{column}
\end{columns}
\end{frame}

\begin{frame}[fragile]{TP B4 --- courbe PR et seuil métier (30 min)}
\begin{codebox}[B4. Choisir le seuil sous contrainte de budget NOC]
\begin{lstlisting}
from sklearn.metrics import precision_recall_curve
prec, rec, thr = precision_recall_curve(yte, proba)

# budget : 5 alertes/jour/serveur
# test = 9 jours x 8 serveurs -> au plus 5*8*9 = 360 alertes
def alertes_jour_serveur(s):
    n = (proba >= s).sum()          # alertes sur tout le test
    return n / 9 / 8                # par jour et par serveur

for s in [0.5, 0.2, 0.1, 0.05, 0.02, 0.01]:
    p = (proba >= s).astype(int)
    print(s, alertes_jour_serveur(s),
          precision_score(yte, p, zero_division=0),
          recall_score(yte, p, zero_division=0))
\end{lstlisting}
\end{codebox}
\begin{abox}[Livrable B4]
La courbe précision-rappel de VOTRE modèle + le tableau
seuil $\rightarrow$ (alertes/jour/serveur, précision, rappel), et le
\textbf{seuil recommandé} qui respecte les 5 alertes/jour/serveur.
Référence du corrigé (gradient boosting) : \textbf{40 alertes sur
9 jours} $= 0{,}56$/jour/serveur, précision 1,00, rappel 0,625.
\end{abox}
\end{frame}

\begin{frame}{TP B5 --- livrable et barème}
\begin{columns}[t]
\begin{column}{0.52\textwidth}
\begin{abox}[À rendre]
\begin{enumerate}
\item \textbf{Tableau comparatif} des modèles (au moins
logistique, forêt, boosting --- et IsolationForest en repère) :
précision, rappel, $F_1$, AP, alertes/jour/serveur au seuil choisi.
\item \textbf{Seuil recommandé} avec sa justification métier
(budget NOC, coût FP vs FN).
\item \textbf{Un paragraphe} (5--8 lignes) : quel modèle vous
déployez, pourquoi, et ce que vous n'avez PAS détecté (les 24 FN
du corrigé --- que farez-vous ?).
\item Reproductibilité : graine 5534, split temporel cité.
\end{enumerate}
\end{abox}
\end{column}
\begin{column}{0.44\textwidth}
\begin{kbox}[Barème (20 pts)]
\begin{itemize}
\item B1 IsolationForest sans fuite --- 3 pts
\item B2 recouvrement + commentaire honnête --- 3 pts
\item B3 forêt balanced + matrice juste --- \textbf{5 pts}
\item B4 courbe PR + seuil sous budget --- \textbf{5 pts}
\item B5 comparatif + argumentation métier --- 3 pts
\item Reproductibilité (graine, split) --- 1 pt
\end{itemize}
\end{kbox}
\begin{wbox}[Fuite = 0 au critère B3]
Toute valeur calculée sur train+test mélangés (standardisation,
seuil, SMOTE) annule le critère : le futur ne sert qu'à juger.
\end{wbox}
\end{column}
\end{columns}
\end{frame}
